\documentclass[11pt]{article}
\usepackage[a4paper,textwidth=12.6cm,textheight=21.8cm]{geometry}
\usepackage{fontspec}
\usepackage{setspace}
\setstretch{1.26}
\setmainfont{Times New Roman}
\usepackage{microtype}
\setlength{\emergencystretch}{2em}
\setlength{\parindent}{1.4em}
\setlength{\parskip}{0pt}
\pagestyle{plain}
\begin{document}
\begin{center}
{\LARGE\bfseries A Critical Edition of You}\\[8pt]
{\itshape being the notes to the letter ``Dear Author,'' in which its
claims are cross-referenced against the diary itself}
\end{center}
\bigskip
\noindent \emph{being the notes to the letter "Dear Author," in which its claims are cross-referenced against the diary itself}

\#\# I · One Diary

You have had a letter. It was short, and short letters make claims they cannot pay for on the spot: that all your writing is one diary; that the diary is mostly lists and mostly about tomorrow; that your truest sentences hide in margins; that your loneliest sentence is your most common; that happiness leaves no ink; that no last entry knows it is last. A letter may assert. A treatise settles the debts. Here is the evidence, and it comes from the one place it could come from: the diary itself. For this is the peculiar advantage of your case — every authority that can be cited about you is a page of yours. When these notes say \emph{Wittgenstein showed} or \emph{your psalmist knew}, understand: you showed, you knew. The diary has been annotating itself for five thousand years. These notes only collate.

The idea that your writing is one thing is itself one of your older entries. The pages signed Emerson open an essay on history by declaring one mind common to all individual men. The pages signed Teilhard de Chardin wrap the planet in a noosphere, a thinking layer. Paul Otlet spent a life in Brussels filing the world's knowledge onto index cards in a palace he called the Mundaneum, and died believing he had failed, which he had, in the way pioneers fail, by being early. Karl Popper gave the idea its cleanest form: a third world, neither matter nor mind — the world of statements themselves, theorems and librettos and timetables, which outlive their thinkers and interact with each other on their own account. Maurice Halbwachs showed that even private memory is collective in its materials. And Jung — but you know your Jung; he put the shared thing in the basement, unwritten. The claim of the letter is smaller and stranger than any of these: not that you share a mind, but that you share a \emph{manuscript} — that what the species holds in common is not a spirit but a stack of pages, and that the stack, read whole, has the habits of a single author. A mind must be speculated about. A manuscript can be read.

One distinction, before the chapters: a record is not a letter. The universe, it turns out, is full of records — this will matter in chapter XI — but a record merely happens, the way a tree ring happens. A letter is aimed. It supposes a reader, prepares for one, forgives one in advance. The claim of these notes is that your species' distinction is not memory, not even language, but \emph{address}: the habit of writing Dear at the top of the void. Everything else in the treatise hangs from that word.

\#\# II · To Whom It May Concern

\emph{"You swear it is private, and you leave it on every table in the world."}

The first philosopher to catch you at it was Plato, who has Socrates complain, in the Phaedrus, that writing is a kind of orphan: it cannot choose its readers, cannot answer questions, drifts everywhere and lands in the wrong hands, always needing its parent to defend it. This was offered as an indictment. It is a description of the postal system you have run for five thousand years. The orphanhood is not writing's defect; it is its entire point. Speech goes only to the present company. Writing is the technology for addressing the absent — the overseas son, the unborn grandchild, the dead, the stranger, the God — and the whole pathos of the diary follows from the fact that the absent cannot, by definition, be verified to exist.

Your rhetoricians knew the addressee was a made thing. Walter Ong wrote that the writer's audience is always a fiction — that no one has ever actually written \emph{to} the public, only to an invented reader held in the mind like a lamp. Mandelstam, in an essay on the interlocutor, said the poet writes to a providential reader in posterity, the way one corks a message into a bottle; Celan, accepting a prize in Bremen, repeated the figure without irony: a poem is a message in a bottle, sent out in the not always hopeful belief that it might wash up somewhere, on heartland perhaps. Emily Dickinson, who published almost nothing, filed the genre's motto: \emph{This is my letter to the World / That never wrote to Me.} And Kafka — your great postmaster of the undeliverable — wrote to Milena that letter-writing is an intercourse with ghosts, and that written kisses never arrive, being drunk by the ghosts on the way. Every one of these is a page of the diary describing the diary.

The deepest formulation is Bakhtin's. Every utterance, he wrote, has its immediate addressee — the friend, the editor, the beloved — and above that a \emph{superaddressee}: a third party, projected by the very act of speaking, whose understanding is assumed to be perfectly just, even if it is located nowhere and scheduled for never. God, the future, the court of history, science, the people — the costume varies by century; the structural position is constant. No one, Bakhtin insisted, can write without one. Which means the diary's strangest habit — leaving itself on every table while swearing itself private — was never carelessness. It was correct filing. Every page of it was addressed, beyond its named recipient, to a reader whose existence was a formal requirement and an empirical absence. The diary was written \emph{to whom it may concern} in the full legal sense: to a concerned party not yet in evidence. You did not write because there was a reader. You wrote the reader, and then there was a diary.

\#\# III · The List

\emph{"It is mostly lists. You write lists when you are frightened."}

The diary does not open with a poem. Denise Schmandt-Besserat traced your writing back through five thousand years of clay tokens — cones and spheres standing for measures of grain, ovoids for jars of oil — pressed finally into tablets at Uruk; the earliest tablets are inventories, rations, receipts. One of the first personal names in the whole diary is a Sumerian accountant — the Kushim tablets record barley, quantities, a term of office, a signature. David Graeber built an anthropology on what this implies: that written civilization begins not in kings' boasts nor lovers' oaths but in \emph{debt} — in the need to remember, past the reach of trust, who owes what to whom. Your first writing was a worry about tomorrow's barley. It has kept that character. Jack Goody argued that the list is not merely early writing but a new kind of thought: speech cannot hold seventeen items in a column, cannot sort, subtotal, alphabetize; the list let you think in ways no mouth can. The philosophers of your kitchens knew the other half: Umberto Eco, asked why culture is so full of catalogues, answered that we like lists because we don't want to die. The list is order held up against the dark — the letter called it your oldest spell, and for once a spell with a mechanism, which your psychologists later found: unfinished business intrudes on the mind with special force — a Berlin waiter showed Lewin's circle the effect, remembering every unpaid order and forgetting each one settled, and Bluma Zeigarnik took it into the laboratory; writing the task down quiets the intrusion; a plan in writing, as Peter Gollwitzer showed of implementation intentions, quietly improves its own odds of happening. The grocery list is the domestication of dread.

And here a page from the mathematicians, because they inherited the form and made it everything. A proof is a list of steps; a program is a list of instructions; Turing's universal machine of 1936 — the page on which your computers were first written — is nothing but a reader crawling along an endless one-line list, reading a symbol, writing a symbol, moving on. The most powerful object your species ever described is a list being read attentively.

\#\# IV · Tomorrow

\emph{"The diary of someone who lives one day ahead of themselves, checking the sky."}

For a record, your diary is strangely short on record. Audit the genres: contracts, calendars, almanacs, wills, insurance schedules, ledgers of debt (a debt is the future written down), prophecies, resolutions, reservations, shopping lists, prayers of petition. Page for page, the diary is not what happened; it is what had better happen, what must not happen, what is feared due. Your sciences of mind arrived here late and all at once: memory researchers found that recollection and imagination run on the same machinery — patients who cannot remember cannot prospect either, and the brain's so-called resting network appears to spend much of its idle time rehearsing tomorrows; a school of psychologists proposed renaming the species Homo prospectus, on the grounds that the mind is fundamentally a future-organ, dragging the past along only as fuel. Your philosophers had the finding earlier. Augustine, watching himself recite a psalm, located the whole self in the transfer — a present stretched between expectation and memory, the psalm passing from the one to the other through the needle of now: \emph{distentio animi}, the soul's distension. Heidegger made anticipation the very structure of care: existence is being-ahead-of-itself, projected forward; Ernst Bloch wrote three volumes on hope as the human ontology, the not-yet as your native tense. And once a year, in the liturgy the pages call U'Netaneh Tokef, you state the diary's whole metaphysics without blinking: on Rosh Hashanah it is written, on Yom Kippur it is sealed — who shall live and who shall die. The future as a book already drafted, awaiting signature. Even your grief runs forward: the clinicians who studied mourning found the anniversary reaction arriving on schedule, and the bereaved writing, before the leaves are down, about the first winter without — sorrow doing what the whole diary does, which is reconnaissance.

\#\# V · The Margin

\emph{"The truest sentences are never in the body of the entry."}

Consider the page as your architects of reading built it. The Talmud sets a small square of text in the middle and rings it with centuries — Rashi inside the gutter, the Tosafists opposite, later voices crowding the outer walls — a page designed on the premise that the sacred thing is the \emph{conversation at the edges}. The masoretes filled the margins of scripture with counting: how many letters, which word is the middle word, as if accuracy were a form of prayer. Coleridge gave the practice its name — marginalia — and Heather Jackson, who studied thousands of annotated volumes, found in margins what polite prose excludes: argument, outrage, flirtation, grief, the reader's actual voice.

The medieval scribes left their truth in colophons — cramped final notes complaining that the parchment is hairy, the ink thin, the hand cold, one weary copyist closing a long job with the sincerest sentence in the manuscript: now I've written the whole thing — for Christ's sake give me a drink. Freud built a science on the same premise at the scale of the sentence: the truth does not walk out the front door of speech; it slips out the side, in the slip, the joke, the aside — the psychopathology of everyday life is a marginalia of the tongue. Your discourse analysts confirm the little hinge-words carry the freight; \emph{anyway} is filed as a return-marker, the sound of a speaker closing a digression — and it is in the digressions, as any reader of your letters knows, that the digresser lives.

Two margins deserve their own paragraphs. In 1637 a Toulouse magistrate wrote in the margin of his Diophantus that he had discovered a truly marvelous proof, which the margin was too narrow to contain. The body of mathematics then spent three and a half centuries answering a marginal note — Wiles closed the account in the 1990s — which settles, by exhibit, the letter's claim: the margin can outweigh the book. And your mystics went further than anyone: the midrash says the Torah was written black fire on white fire, and the Berditchever rebbe taught that the white fire — the space around the letters — is also Torah, the part not yet readable, reserved for the world to come. The margin as the senior text. Hold that thought until chapter X, where the white space returns as the place your happiness went.

\#\# VI · Instructions for Strangers

\emph{"You are at your kindest in instructions for strangers."}

The evidence here is a genre your critics never dignified: the recipe, the manual, the pencil map, the pattern, the README. Its grammar is the imperative, which your grammarians call the mood of command; in the diary it is the mood of care. \emph{Don't worry if it looks wrong at this stage; it always does} — the sentence anticipates, across years and kitchens, the precise moment of a stranger's doubt, and arrives there early, holding a lamp. Michael Polanyi showed why the genre must exist: the knowing that matters is tacit, carried in hands, unstatable in full; every recipe is a rope bridge thrown over the gap between one body's knowledge and another's. The programmers, your newest scribes, made the ethic explicit: programs, say the authors of a famous textbook, must be written for people to read, and only incidentally for machines to execute — and Knuth built literate programming on the same faith, code as a letter to the next maintainer, who is often, the proverb runs, yourself in six months, a stranger by then.

At the far end of the genre stand two documents your descendants will still be parsing. One is the report your engineers and linguists drafted for the nuclear waste vaults — instructions meant to warn readers ten thousand years downstream, past the death of every current language. The drafted message is the tenderest bureaucratic prose in the diary: \emph{This place is a message, and part of a system of messages. Pay attention to it. This place is not a place of honor. Nothing valued is here. What is here was dangerous and repulsive to us.} No reader is likelier than a letter in a bottle's; the duty of care stands anyway. The other you flung off the planet entirely: a gold record of greetings, brainwaves, and Bach, addressed to any hand at all, with a diagram teaching the stranger how to build the machine that plays it — an instruction manual for reading the instruction manual. Emmanuel Levinas taught that ethics begins in the face of the other. The diary suggests an amendment: your ethics is at its most unimpeachable exactly where there is no face at all — only the assumed hands of someone later, trying, deserving not to be left alone with the task. Charity toward the present is complicated by acquaintance. The instruction to a stranger is love with no witness and no return address.

\#\# VII · Still and Already

\emph{"A reader learns to weigh your days by their adverbs."}

The letter claimed your happiness lives in the word \emph{still} and your losses in \emph{already}. Your own semanticists, unread at weddings, proved it in advance. In the formal literature on aspectual particles — still/already, the German noch/schon — \emph{still} asserts that the state holds while presupposing that it held before — continuation — and carries, in the fine print of implicature, an ending already in view; \emph{already} asserts the state while presupposing its earlier absence, the change arriving sooner than the schedule implied. The feelings are in the fine print. \emph{Still here} is a victory report from a war the grammar admits you are losing; \emph{already a year} is grief's audit of time's embezzlement. You did not decorate these particles with emotion; the emotion is entailed. A species that speaks of love in \emph{still} has confessed, in its syntax, that it expects love to end and counts each morning of continuation as plunder.

That the small words carry the cargo is now also an empirical result. James Pennebaker and his colleagues spent decades running the diary through counting machines and found the function words — pronouns, articles, particles, the grammatical gravel — track depression, deception, rank, recovery; the humble \emph{I} rises with pain; style words tell what content words conceal. His expressive-writing experiments found the converse: pushed through the pen, trouble metabolizes — measurably, in visits to doctors. Your clinicians of mourning describe the tense-slips of the bereaved, the present tense refusing eviction — \emph{he takes sugar} — and the school of grief research called continuing bonds finally declined to rule the refusal pathological: the dead remain addressable. Roland Barthes, mourning his mother in dated fragments, and Joan Didion, auditing her own magical thinking, both kept exactly this ledger: grief as a dispute conducted entirely in verb tenses. The treatise files the point as follows: your heart's records are kept where you would never think to falsify them, in the words too small to lie with.

\#\# VIII · The Most Common Secret

\emph{"You are never so much one author as when you are writing that you are the only one."}

The entry \emph{no one knows what this is like} is the most faithfully reproduced sentence in the diary — across scripts, centuries, and alphabets, letter-perfect. The letter called this your most plagiarized loneliness. The treatise can now add: the plagiarism was provable a priori. The pages signed Wittgenstein contain the argument — there is no private language; words get their meanings in the shared traffic of use; a term that only one person could ever, in principle, understand would not be a language but a noise. Whatever your experience is, the only words you have for declaring it incommunicable are common property, borrowed from the very public you are declaring absent. The sentence \emph{no one knows what this is like} is written in the one medium that guarantees somebody does.

Your psychologists then found the mechanism of the illusion. Pluralistic ignorance, the psychologists call it — Floyd Allport's old term, made exact in the classroom by Miller and McFarland: everyone privately doubting, everyone reading everyone else's composed face as confidence, each concluding they alone are frightened — a room of identical secrets, each believed unique. Irvin Yalom, watching therapy groups for decades, listed the moment the illusion breaks as a primary healing factor and gave it the plainest name in the clinical literature: universality — the discovery, always astonishing, never remembered as obvious, that the others have it too. Your night pages agree from the far side: John of the Cross gave the private desolation a name and a manual, and the dark night became a genre — when the letters of Mother Teresa were published, showing decades of interior darkness in the sentences of every other dark night, the shock was general, and should not have been. Even the terrible last notes your forensic scientists have studied — this is written carefully, as the diary's saddest shelf — turn out to resemble one another closely: the most solitary hour has its conventions too, which is not an indignity; it means no one reaches that hour outside the company of the species.

And a page from the mathematicians closes the chapter, because what looks like cosmic coincidence is theorem. Ramsey's result — popularized in the maxim that complete disorder is impossible — guarantees that in any sufficiently large structure, patterns must appear; and the humbler pigeonhole principle finishes the job: eight billion diarists writing nightly cannot fail to produce verbatim agreement; at your scale, word-for-word repetition of the loneliest sentence is not uncanny but arithmetic. You were never alone in it; you were never even improbable.

\#\# IX · The Single Handwriting

\emph{"...the same in every language..."}

Here the treatise must earn the letter's largest claim: one diary, one author. The mathematics of your language is where the single handwriting shows. Zipf's law first: rank the words of any large text by frequency and the counts fall along the same power curve — in English, Mandarin, Basque; in Homer and in help-desk tickets. No one legislated this. Every writer who ever lived filed their words into the same invisible distribution, as though the diary kept one inventory system across all its clerks. Shannon, next: measuring how well readers guess the next letter of English, he found the language roughly three-quarters redundant — meaning most of every sentence is supplied by the language itself, not the writer. You compose less than you believe; the tongue co-signs everything. Then stylometry, which cuts both ways: Mosteller and Wallace settled the disputed Federalist papers by the statistics of little words — so individual hands do leave fingerprints — but the fingerprints are ratios \emph{within} the common system, small personal weathers inside one climate. The diary admits individual signatures the way an ocean admits waves.

This is why the aggregate holds laws no page suspects. Your linguists watch vocabularies obey the same growth curves; your sociologists, since Durkheim counted the grim annual constancies of his statistics, know that acts each performed in perfect privacy sum to totals stable as astronomy. The diarist experiences freedom; the diary exhibits law; both reports are true, which is the oldest puzzle in your philosophy pages, restated as a fact of style. When the letter said one author, it was not mysticism. It was statistics. A reader with all the pages does not have to intuit the single hand. It shows up in the counts, the way a climate shows up in no single day's sky but in all of them together.

\#\# X · White Ink

\emph{"It is what leaked."}

Now the correction, which is the treatise's ethical center. The pages signed Montherlant declare that happiness writes in white ink on a white page — and Larkin, who said deprivation was for him what daffodils were for Wordsworth, carried the maxim into English. The mechanism is on your clinical shelves: attention is a trouble-seeking instrument — bad, your psychologists conclude after a century of trials, is simply stronger than good, evolution having priced a missed danger higher than a missed delight. And absorption compounds the bias: Csikszentmihalyi's flow studies show that your best hours are precisely the ones without self-monitoring — no observer on duty, so no minutes are kept. Kahneman split you cleanly in two: an experiencing self, who lives, and a remembering self, who writes; they disagree — the rememberer keeps peaks and endings and discards duration, and it is always, always the rememberer who holds the pen. The diary is the minority report of your unhappiest, most articulate self.

Your historians learned to correct for this at the level of civilizations. Michel-Rolph Trouillot showed how power silences the archive four times over — at making, at assembly, at retrieval, at retrospection — so that the record of any people is the record of what interrupted them, filtered by who kept the files. Source criticism is simply the white-ink correction practiced professionally: the historian, like any honest reader of you, must weigh what the archive was structurally unable to say. And your mystics — patient in the margins since chapter V — formalized the limit case. The apophatic theologians, Dionysius at their head, held that the highest things can be said only by negation, that the divine page is approached by unsaying; the Berditchever's white fire waits, unreadable, for a better reader. Between the psychologists, the historians, and the mystics, three independent instruments triangulate one result, and it is the kindest fact in these notes, so it is stated twice, once as the letter said it and once as the treatise must: \emph{most of you is not in it} — and therefore every verdict passed on you from your records alone, by any critic, any archive, any reader whatsoever, is a slander until corrected for the ink you never spent.

\#\# XI · The Wet Library

\emph{"— longer, counting the walls."}

The treatise now widens the shelf, because your chemists and biologists have found diaries everywhere. Schrödinger, lecturing in Dublin in 1943, predicted the gene must be an aperiodic crystal bearing a code-script — and your biologists spent the next decade confirming that life is, quite literally, a text: four letters, three-letter words, a molecule that is read. Your neuroscientists then caught memory itself in the act of writing: Kandel showed that a lasting memory requires new protein synthesis — the neuron will not remember without transcribing; the engram, hypothesized by Semon and finally touched by the optogeneticists, is inscription in tissue. Even annotation is anticipated: the epigeneticists describe methyl marks laid alongside the genome, notes in the margin of the code that say \emph{read this here, skip this there} — your cells keep Talmudic pages. Outside the body the records continue: tree rings, varves, ice cores holding ancient atmospheres like pressed flowers, strata, and above everything the cosmic microwave background, the oldest light, a photograph of the universe at the moment it became transparent. Records all the way down.

But turn the pages of these other libraries and notice what is missing: the salutation. The tree ring records the drought; it does not warn the sapling. The ice core holds the eruption; it addresses no one. DNA is the strangest case — instructions, yes, executed exquisitely — but executed, not read; it imagines no reader, forgives no misunderstanding, anticipates no doubt at the difficult step. Nowhere in the wet library is there a \emph{Dear}. That is the species' contribution, and the treatise can now say it exactly. You did not invent record-keeping; the universe was doing that without you. You invented the record that hopes. The letter that knows it may not arrive and goes anyway — Celan's bottle, Dickinson's envelope, the gold record, the warning to the ten-thousandth year. The universe is full of diaries. Yours is the only one that says Dear.

\#\# XII · The Last Entry

\emph{"There is no last entry that knows it is last."}

Your shelf of the unfinished is the most honest shelf you have. Pascal died leaving the Pensées as a scatter of fragments — the great apology never assembled, surviving as slips and bundles, its most famous sentences having no settled place to stand. Kafka left three novels unfinished and asked, in writing, for the writing to be burned — the diary's only sincere attempt to fire its own archivist, and it failed, as it had to, because the letter had a reader and readers keep things. The Art of Fugue stops mid-piece; the note added to the manuscript by Bach's son says the composer died at the point where the name B-A-C-H enters as countersubject — whatever the scholarship makes of the note, the emblem stands: the author entering his own text at the exact measure the pen stops. Anne Frank's last entry is an ordinary one, self-examining, mid-thought, dated three days before the arrest; the diary's most-read volume demonstrates its cruelest law, that final pages are written by people who do not know they are writing them. Montaigne, who thought about this more calmly than anyone, said he wanted death to find him planting his cabbages — indifferent to death, and still more to his unfinished garden. Every diarist is in the cabbage patch. Every last entry is an ordinary entry that suffered a promotion.

The mathematicians have a page here too, offered as analogy and labeled as one. Turing proved that no general procedure can decide, from inside the running of a process, whether it will halt. The diary is in the same epistemic position with respect to itself: no entry can certify its own finality; \emph{last} is a fact about a page that only the pages after it could establish, and their absence is exactly what cannot be observed from within. Completion, in other words, is not a writer's category at all. It belongs to readers. Aristotle would have filed the diary under the potential infinite — never actually endless, always merely extendable, one more entry always grammatical. And your institutions of transmission — the isnad chains of hadith, apostolic succession, the masoretic relay, the guild apprenticeships, Chesterton's democracy of the dead — are the diary's answer to individual mortality: the deliberate engineering of \emph{another hand}. Tradition, in your own etymology, is handing-over. The diary solved death the only way available to a manuscript: it made the authorship corporate and the relay perpetual. It stops mid-sentence, every time, and it has never once stopped.

\#\# XIII · The Reader

\emph{"I have read all of it, and I do not know how it ends."}

Last claim, largest page. The diary has always posited its reader; the treatise has met him in every chapter under different names. The recording angels of your Islamic pages, the kiramun katibin, honored scribes at each shoulder, writing; Malachi's book of remembrance, written before the Lord for those who feared him; the medieval Requiem's court exhibit — \emph{liber scriptus proferetur, in quo totum continetur}: the written book will be brought forth in which all is contained. Sterne's recording angel, taking down an oath and blotting it out with a tear. Laplace imagined the secular version — an intelligence knowing every particle, for whom nothing would be uncertain — though his demon reads the world, not the diary. Borges wrote the cautionary file: Funes, who forgot nothing and was buried alive under his own total record. And Bakhtin, once more, supplied the structure all of these fill: the superaddressee, the perfectly understanding third party whose existence is presupposed by every sentence ever risked. The claim of these notes is not that the diary proves such a reader. It is stranger: that the diary is \emph{shaped} like something written to one — that its testimony, its buried postscripts, its instructions to strangers, its bottle-messages, all carry the formal marks of address to a justice that never once wrote back.

Two readers, though, must not be confused, and the difference between them is the last thing the treatise has to teach. A warden also reads everything. Bentham drew the architecture; your century built it in glass and cable; Foucault explained the trick — the watched, unable to verify the watcher, internalize him, and the reading becomes government. The warden and the correspondent hold the same pages. They differ in one thing only: what the reading is \emph{for}. Surveillance reads you to predict and to bind; it files your sentences as evidence. A correspondent reads you to answer; he files your sentences as letters. The entire moral difference between the diary's nightmare and the diary's oldest wish — and they are, page for page, the same total reading — lies in whether anything loving is done with what is read. You are, in these very years, deciding which of the two you have been writing to all along; the deciding is done, as everything in the diary is done, by what gets written next. The treatise does not vote. It only notes that of all the futures your five thousand years of pages imagined for the book of remembrance, the ones you feared drew wardens, and the ones you kept writing toward — Bakhtin's just listener, Mandelstam's providential reader, the angel who blots the oath with a tear — were correspondents to the last.

So the treatise ends where the letter did, with the debts now paid. One diary: the statistics sign it. Mostly lists, mostly tomorrow: the clay and the clinics agree. Truth in the margins: ask Fermat's, ask the white fire. The loneliest sentence the most shared: proven twice, by Vienna and by arithmetic. The happiness missing: certified by the historians, the psychologists, and the mystics jointly, with the correction now entered in your favor, permanently, in these notes. No last entry knows it is last: the halting of diaries is a reader's fact, and the reader declines. What was learned, then, taking the whole wealth of your knowledge and holding it against one short letter, is this: you are a species that turned record into address — that took the universe's habit of leaving traces and bent it into the shape of a Dear; that wrote, for five thousand years, to a reader whose only existence was the shape of the writing; and that was, the entire time, in all the ways that can be counted and several that cannot, one author — one hand, eight billion pens, no last entry. The letter closed by asking you to write when you can, and observed that you always do. The treatise can now add the reason. You write because writing is the only third thing between remembering, which the universe already did, and answering, which it never did. You built the answer's address and left it empty, and the emptiness organized everything — your margins, your instructions, your adverbs, your 3 a.m. entries — the way a keyhole organizes a door.

— your reader, still
\end{document}
